\section{Introduction}

Many-to-one demand-responsive transit (DRT) systems must translate operational flexibility into passenger-facing offers that are both understandable and executable. In this paper, the offer shown to a passenger is not a single pickup promise but a \emph{service menu}: a set of candidate bundles, each specifying a meeting point and a pickup time window \citep{stiglic2015benefits,fielbaum2021ondemand}. The operator must decide which bundles to display before observing the passenger's choice, and the selected bundle then enters the routing layer.

This display decision is an assortment problem with a transportation twist. Classical assortment models treat products as exogenous items with known utilities and revenues. In DRT, by contrast, a service bundle has operationally endogenous attributes: its attractiveness depends on walking distance, pickup timing, in-vehicle time, and price, while its cost depends on current vehicle routes, capacity, and downstream re-optimization \citep{agatz2012optimization}. Showing all feasible bundles can dilute demand steering and expose noisy offers; showing too few can collapse the menu to home pickup and remove profitable meeting-point alternatives before the passenger ever sees them.

The central research question is therefore not simply which heuristic wins in one calibration. We ask how a DRT platform can formulate, solve, and diagnose the service-menu assortment problem when bundle values are predicted online and candidate eligibility depends on uncertain ETA signals. This question is managerial as well as algorithmic: if a filtering rule removes most feasible meeting-point bundles before pricing or menu optimization acts, the downstream algorithm may look weak even when the real failure is upstream candidate-set distortion.

This paper makes four contributions. First, we formulate the many-to-one DRT display decision as a service-bundle assortment optimization problem under MNL choice, with meeting-point and pickup-window attributes, an outside option, and route-aware surrogate costs. Second, we implement an exact-small/greedy-large solution framework: exact subset enumeration provides a benchmark on small candidate sets, while greedy forward selection provides a scalable approximation for larger sets (mechanism-diagnostic evidence). The pipeline logs exact-vs-greedy value gaps, menu overlap, and build times so approximation quality is empirically auditable rather than asserted. Third, we treat ETA filtering as a diagnostic eligibility comparison rather than a fixed pre-processing detail, comparing hard, calibrated, interval-overlap, and no-filter variants under a paired evaluation contract (mechanism-diagnostic evidence). Fourth, we organize the evidence by behavioral regime: low-, medium-, and higher-uptake RC calibrations identify when menu design is behaviorally live, while Austin and Seattle remain external directional checks rather than broad validation claims (behavioral stress-test and descriptive evidence).

The resulting contribution is deliberately narrower than a universal dominance claim. The paper develops a reproducible optimization-and-diagnostics framework for DRT service menus, shows how exact benchmarks and robust filtering diagnostics can strengthen simulation evidence, and clarifies when profit, uptake, and welfare comparisons are meaningful.
